% Source-derived chapter 04: Shift Operations
% Original source: rigid-local-systems-multiplicative-eigenvalue-problem
\chapter{Shift Operations}
\label{rigid-local-systems-multiplicative-eigenvalue-problem-chapter-04}
\source{rigid-local-systems-multiplicative-eigenvalue-problem}

\begin{remark}[Source section]
\label{rigid-local-systems-multiplicative-eigenvalue-problem-chapter-04-source}
\source{rigid-local-systems-multiplicative-eigenvalue-problem}
\source{rigid-local-systems-multiplicative-eigenvalue-problem}
This chapter reproduces the corresponding section of the retrieved
LaTeX source, including its definitions, statements, examples, and
proofs. It has not yet been translated into Lean.
\notready
\end{remark}

% The source section title was: Shift Operations
\label{shifty}
\source{rigid-local-systems-multiplicative-eigenvalue-problem}
Let $p=p_i\in S$. We define a shift operation $\operatorname{Sh}_p: \Par_{n,\mn,S}\to \Par_{n,\mn(p),S}$,  as follows:
$$\operatorname{Sh}_p(\mw,\me,\gamma)=(\mw',\me',\gamma')$$
where
\begin{enumerate}
\item $\mw'$ coincides with $\mw$ away from $p$, and is defined in a neighbourhood of $p$ as follows: Let $t$ be a uniformizing parameter at $p$, $U$ an open subset containing $p$, and assume $t$ does not vanish on $U-\{p\}$. Sections $s$ of $\mw'$ on $U$ are meromorphic sections of $\mw$ of $U$, regular on $U-\{p\}$ such that $ts$ is a regular section of $\mw$
whose fiber at $p$ is in $E^p_1$. It is easy to see that $\mw'$ is independent of the choice of the uniformizing parameter $t$ at $p$. Clearly $\mw\subseteq\mw'$.
\item
The natural map $\det(\mw)\to \det(\mw')$ has a zero of order 1 at $p$, and hence $\det(\mw')$ is isomorphic to $\mn(p)$. Let $\gamma'$ denote this isomorphism.
\item Since $\mw$ and $\mw'$ are isomorphic outside of $p$, we can take $E'^{q}_{\bull}=E^q_{\bull}$ for $q\in S-\{p\}$. To complete the definition of $\me'$ we need to define $E'^p_{\bull}$.  Now $\mw_p\to \mw'_p$ has kernel $E^p_1$, so we define $E'^p_a$ to be the image of $E^p_{a+1}$ for $a<n$ and equal to  $\mw'_p$ for $a=n$.
\end{enumerate}

The exact sequence
$$0\to \mw \to \mw'\to \mw'_p/E'^p_{n-1}\to 0$$
allows us to define the inverse $\operatorname{ISh}_p$ to $\operatorname{Sh}_p$.
\subsection{The effect of the shift operation on line bundles}
Writing
$\operatorname{Sh}_p(\mw,\me,\gamma)=(\mw',\me',\gamma')$
we have an exact sequence
\begin{equation}\label{roast1}
\source{rigid-local-systems-multiplicative-eigenvalue-problem}
0\to \mw'\to \mw(p)\leto{t} E^p_n/E^p_1\to 0
\end{equation}
which we make independent of the choice of $t$ by twisting by the canonical bundle:
\begin{equation}\label{roast}
\source{rigid-local-systems-multiplicative-eigenvalue-problem}
0\to \mw'\to \mw(p)\to(E^p_n/E^p_1)\tensor K_p^*\to 0
\end{equation}
\begin{lemma}
\source{rigid-local-systems-multiplicative-eigenvalue-problem}
\notready\label{phase}
\source{rigid-local-systems-multiplicative-eigenvalue-problem}
$E^p_a/E^p_{a-1}=E'^p_{a-1}/E'^p_{a-2}$ if $2\leq a\leq n$, and $E^p_1/E^p_0=(E'^p_n/E'^p_{n-1})\tensor K_p$.
\end{lemma}

\begin{proposition}
\source{rigid-local-systems-multiplicative-eigenvalue-problem}
\notready\label{easily}
\source{rigid-local-systems-multiplicative-eigenvalue-problem}
Let $p=p_i\in S$.
\begin{enumerate}
\item[(a)] $\operatorname{Sh}_p^*(\mathcal{B}(\vec{\lambda},\ell))= \mathcal{B}(\vec{\mu},\ell)$
where $\mu^j=\lambda^j$ for $j\neq i$ and $\mu^i_a= \lambda^i_{a-1}$ if $2\leq a\leq n$, and $\mu^i_1=\lambda^i_{n}+\ell$.
\item[(b)] $\operatorname{ISh}_p^*(\mathcal{B}(\vec{\lambda},\ell))= \mathcal{B}(\vec{\nu},\ell)$
where $\nu^j=\lambda^j$ for $j\neq i$ and $\nu^i_a= \lambda^i_{a+1}$ if $1\leq a< n$, and $\nu^i_{n}=\lambda^i_{1}-\ell$.

\item[(c)]The morphism $\operatorname{Sh}_p: \Par_{n,\mn,S}\to \Par_{n,\mn(p),S}$ induces bijections $\Pic^+(\Par_{n,\mn(p),S})\leto{\sim}\Pic^+(\Par_{n,\mn,S})$ and
$\Pic^+_{\Bbb{Q}}(\Par_{n,\mn(p),S})\leto{\sim}\Pic^+_{\Bbb{Q}}(\Par_{n,\mn,S})$.
\item[(d)] $\ml\in  \Par_{n,\mn(p),S}$ is an F-line bundle if and only if $\operatorname{Sh}_p^*\ml$ is an F-line bundle on  $\Par_{n,\mn,S}$.
\end{enumerate}
\end{proposition}
\begin{proof}
$\operatorname{Sh}_p$ and $\operatorname{ISh}_p$ are isomorphisms (they are inverses of each other). This proves (c). Lemma \ref{phase}, and \eqref{roast1}  (and Definition  \ref{levell}) imply parts (a) and (b). \end{proof}

\begin{defi}
\source{rigid-local-systems-multiplicative-eigenvalue-problem}
\notready\label{grade} Define maps: $\operatorname{Gr}:\Pic (\Par_{n,\mn,S})\to \Bbb{Z}/n\Bbb{Z}$ by
\source{rigid-local-systems-multiplicative-eigenvalue-problem}
$(\vec{\lambda},\ell)\mapsto \sum_{i=1}^s|\lambda_i|+ \ell \deg\mn \in\Bbb{Z}/n\Bbb{Z}$. Note that shift operations preserve this grading.  It is also easy to verify that the grade of an effective line bundle is zero (reduce to the case $\deg\mn=0$ using shift operations).
\end{defi}

\begin{remark}
\source{rigid-local-systems-multiplicative-eigenvalue-problem}
\notready\label{able}
\source{rigid-local-systems-multiplicative-eigenvalue-problem}
Recall the Killing form isomorphism $\frh_n^*\to \frh_n$ from Section \ref{loudvoice}. In case (a) of Proposition \ref{easily}, we note that $\kappa(\frac{{\mu^j}}{\ell}),\kappa(\frac{\lambda^j}{\ell})\in \Delta_n$ give the same conjugacy class in $\op{SU}(n)$ if $j\neq i$ and $\kappa(\frac{{\mu^i}}{\ell})$ equals  $\zeta_n^{-1}\cdot \kappa(\frac{{\lambda^i}}{\ell})$ (see Definition  \ref{symmetry}). In case (b),
we have  $\kappa(\frac{{\nu^j}}{\ell})=\kappa(\frac{{\lambda^j}}{\ell})$ if $j\neq i$ and $\kappa(\frac{{\nu^i}}{\ell})=\zeta_n\cdot\kappa(\frac{{\lambda^i}}{\ell})$.
\end{remark}
\begin{lemma}
\source{rigid-local-systems-multiplicative-eigenvalue-problem}
\notready\label{threepointone}
\source{rigid-local-systems-multiplicative-eigenvalue-problem}
Tensoring $\mw$ with $\mathcal{O}(p)$ sets up an isomorphism $\Par_{n,\mn,S}\to \Par_{n,\mn(np),S}$. This isomorphism coincides with an $n$-fold composition of
$\operatorname{Sh}_p$. This isomorphism carries the line bundle $\mathcal{B}(\vec{\lambda},\ell)$ on one stack
to a line bundle with the same parameters $\mathcal{B}(\vec{\lambda},\ell)$ on the other.
\end{lemma}
\subsubsection{Proof of Proposition \ref{LaSo}}\label{shifty2}
\source{rigid-local-systems-multiplicative-eigenvalue-problem}
Proposition \ref{LaSo}, which describes the Picard group of $\Par_{n,\mn,S}$ follows from (a), (b), (c) of Proposition \ref{easily} and the special case $\mn=\mathcal{O}$ proved in \cite{LS}.
\subsection{The effect of shift operations on cycles}
Let $p=p_i\in S$. To study the effect of the shift operations on the stacks $\Omega^0(d,r,\mn,n,\vec{J})$ defined in Definition  \ref{mindful}, we begin with a basic commutative square:
\begin{equation}\label{basicD}
\source{rigid-local-systems-multiplicative-eigenvalue-problem}
\xymatrix{
 \Omega^0(d,r,\mn,n,\vec{J})\ar[d]\ar[r]^{\operatorname{Sh}_p}  & \Omega^0(d',r,\mn(p),n,\vec{K})\ar[d]\\
\Par_{n,\mn,S}\ar[r]^{\operatorname{Sh}_p}  & \Par_{n,\mn(p),S}
}
\end{equation}
where the terms $d'$ and $\vec{K}$ are determined as follows,
 \begin{enumerate}
 \item $d'=d$ if $1\notin J^i$, and $d'=d-1$ if $1\in J^i$
 \item $K^k=J^k$ for $k\neq i$.
 \item $K^i =\{j^i_1-1,j^i_2-1,\dots, j^i_r-1\}$ if $1\notin J^i$, and $K^i =\{j^i_2-1,\dots, j^i_r-1,n\}$ if $1\in J^i$.
 \end{enumerate}

To see that we note that there is a bijective correspondence between subbundles of bundles $\mw$ and $\mw'$ when we have a relation $\Sh_p(\mw,\me,\gamma)=(\mw',\me',\gamma')$. Keeping track of the Schubert cell data gives Diagram \eqref{basicD}.

Recall the cycles $C(d,r,\mn,n,\vec{J})$ on $\Par_{n,\mn,S}$ defined as pushforwards of $\Omega(d,r,\mn,n,\vec{J})$ (Definition  \ref{cyclist}).
 \begin{lemma}
\source{rigid-local-systems-multiplicative-eigenvalue-problem}
\notready\label{newS} Let $p=p_i\in S$. Then
\source{rigid-local-systems-multiplicative-eigenvalue-problem}
 $$\operatorname{Sh}_{p}^*C(d',r,\mn(p),n,\vec{K})=C(d,r,\mn,n,\vec{J}).$$
 \end{lemma}
\begin{proof}
 The operation $\op{Sh}_p$  lifts to a map  $\Omega^0(d,r,\mn,n,\vec{J})\to \Omega^0(d',r,\mn(p),n,\vec{K})$ (as in \eqref{basicD}), but  may not lift to a map of the closures $\Omega(d,r,\mn,n,\vec{J})\to \Omega(d',r,\mn(p),n,\vec{K})$. Since the multiplicities of the divisor
$C(d',r,\mn(p),n,\vec{K})$ can be calculated at generic points, we obtain the stated result from Diagram \eqref{basicD} and Proposition \ref{maya}.
\end{proof}
%By a simple computation the  codimensions of $C(d',r,\mn(p),n,\vec{K},S)$ and $C(d,r,\mn,n,\vec{J},S)$ coincide in the above lemma.
 \subsection{Generalized Gromov-Witten numbers}\label{stuffed}
\source{rigid-local-systems-multiplicative-eigenvalue-problem}
 Definition \ref{GW} is an enumerative count of subbundles on $\mathcal{O}^{\oplus n}$, which is the ``general" rank $n$ vector bundle with trivial determinant. We can replace the role of $\mathcal{O}^{\oplus n}$ by a general rank $n$ vector bundle
with arbitrary determinant $\mathcal{N}$ of degree $-D$, \cite[Definition 2.4]{BHorn}:
\begin{defi}
\source{rigid-local-systems-multiplicative-eigenvalue-problem}
\notready\label{GWgen}
\source{rigid-local-systems-multiplicative-eigenvalue-problem}
Let $I^1,\dots I^s$ be subsets of $[n]$ each of cardinality $r$ and $d\in\Bbb{Z}$. Let $\mn$ be a line bundle on $\Bbb{P}^1$
of degree $-D$.  Let $\mw$ be a general rank $n$ bundle of degree $-D$ (hence $\det\mw=\mn$).  Let $\me\in\Fl_S(\mw)$ be a general point.

The Gromov-Witten number $\langle\sigma_{I^1}, \sigma_{I^2},\dots,\sigma_{I^s}\rangle_{d,D}$
is defined to be the number of subbundles (count as zero if infinite) $\mv\subseteq \mw$ of degree $-d$ and rank $r$  such that for $p\in S$,
$\mv_p\in\Omega_{I^j}(E^p_{\bullet})\subseteq \Gr(r,\mw_p)$. It follows from the dimension formula for quot schemes that
 if $\langle\sigma_{I^1}, \sigma_{I^2},\dots,\sigma_{I^s}\rangle_{d,D}$ is non-zero, then
 \begin{equation}\label{numerical}
\source{rigid-local-systems-multiplicative-eigenvalue-problem}
 \sum_{j=1}^s |\sigma_{I^j}| =r(n-r)+dn-Dr.
 \end{equation}

\end{defi}
These new Gromov-Witten numbers reduce to the usual Gromov-Witten numbers using shift operations by the following:
Assume that the  codimension of the cycle $C(d,r,\mn,n,\vec{J})$ in $\Par_{n,\mn,S}$ is $0$, then in the setting of Lemma \ref{newS}, 
$$\Omega(d,r,\mn,n,\vec{J},S)\to \Par_{n,\mn,S}, \ \ \Omega(d',r,\mn(p),n,\vec{K},S)\to \Par_{n,\mn(p),S},$$
 have the same degree. Let $D=-\deg(\mn)$. We then have by Lemma \ref{newS},  
\begin{equation}\label{formula1}
\source{rigid-local-systems-multiplicative-eigenvalue-problem}
\langle\sigma_{J^1}, \sigma_{J^2},\dots,\sigma_{J^s}\rangle_{d,D}= \langle\sigma_{K^1}, \sigma_{K^2},\dots,\sigma_{K^s}\rangle_{d',D-1}.
\end{equation}
Also note the formula
\begin{equation}\label{formula2}
\source{rigid-local-systems-multiplicative-eigenvalue-problem}
\langle\sigma_{J^1}, \sigma_{J^2},\dots,\sigma_{J^s}\rangle_{d+r,D+n}=\langle\sigma_{J^1}, \sigma_{J^2},\dots,\sigma_{J^s}\rangle_{d,D}
\end{equation}
 which is obtained by an $n$-fold application of the shift operation at any $p_i$.

\subsection{Quantum cohomology of Grassmannians and ranks of conformal blocks}\label{QCB}
\source{rigid-local-systems-multiplicative-eigenvalue-problem}
There is an expression for  $\langle\sigma_{I^1}, \sigma_{I^2},\dots,\sigma_{I^s}\rangle_{d,D}$ (with notation from Definition \ref{GWgen}) as equal to  a suitable  $h^0(\Par_{r,\mathcal{O}(-d),S},\ml)$ \cites{Gepner,Witten} (and Section \ref{QCB} below). The formula in general (even when $D=0$) has some shifting. The most succinct way of stating it is to reduce to $d=0$ (using \eqref{formula1}) and then using the following: Assume $d=0$ in Definition \ref{GWgen} and hence $\sum_{j=1}^s |\sigma_{I^j}|$ is divisible by $r$.
Let $\lambda^1,\dots,\lambda^s$ be highest weights for $\op{SL}(r)$ given by $\lambda^i=\lambda(I^j)$ following Definition \ref{correspondence} below. Then  (see \cite[Proposition 3.4]{BHorn}),
\begin{lemma}
\source{rigid-local-systems-multiplicative-eigenvalue-problem}
\notready\label{newlabel2}
\source{rigid-local-systems-multiplicative-eigenvalue-problem}
Suppose $d=0$ in Definition \ref{GWgen}. Then
$$\langle\sigma_{I^1}, \sigma_{I^2},\dots,\sigma_{I^s}\rangle_{d,D}= h^0(\Par_{r,\mathcal{O},S},\ml)$$ where $\ml=\mathcal{B}(\vec{\lambda},n-r)$.
\end{lemma}
Note that the $h^0$ above is also the rank of a vector space of conformal blocks for the Lie algebra $\mathfrak{sl}_r$ and the representations $\vec{\lambda}$ placed at the points of $S$ at level $n-r$.
\begin{defi}
\source{rigid-local-systems-multiplicative-eigenvalue-problem}
\notready\label{correspondence}
\source{rigid-local-systems-multiplicative-eigenvalue-problem}
Given a subset $I=\{i_1<i_2<\dots<i_r\}\subseteq [n]=\{1,\dots,n\}$, $|I|=r$, let $\lambda(I)=(\lambda_1,\dots,\lambda_r)$ where
$\lambda_a=n-r+a-i_a,\ a=1,\dots,r.$
\end{defi}
If $d\neq 0$, there is a similar equality. Let $\lambda^j=\lambda(I^j)$, and $\ml=\mathcal{B}(\vec{\lambda},n-r)$ on $\Par_{r,\mathcal{O}(-d),S}$. The assumption \eqref{numerical} shows that the grade of $\ml$ is zero (see Definition  \ref{grade}) and using the shift operations and Lemma \ref{newlabel2} we obtain
\begin{equation}\label{newlabel3}
\source{rigid-local-systems-multiplicative-eigenvalue-problem}
\langle\sigma_{I^1}, \sigma_{I^2},\dots,\sigma_{I^s}\rangle_{d,D}= h^0(\Par_{r,\mathcal{O}(-d),S},\ml).
\end{equation}
One can use shift operations to reduce the $h^0$ to the rank of a suitable vector space of conformal blocks.
%The identifications of Proposition \ref{grader} also preserve this grading.

\subsection{Quantum saturation and Fulton conjectures}
We recall quantum saturation \cite{BHorn} and a quantum generalization of a conjecture of Fulton \cite[Remark 8.5]{BKq}. These (now theorems) generalize statements in the classical invariant theory of the special linear group \cite{KT,KTW}.

\begin{proposition}
\source{rigid-local-systems-multiplicative-eigenvalue-problem}
\notready\label{qTheorems}
\source{rigid-local-systems-multiplicative-eigenvalue-problem}
Suppose $\mathcal{B}(\vec{\lambda},\ell)$ is a line bundle on  $\Par_{n,\mn,S}$ of grade zero, where $\lambda^i, i=1,\dots, s$ are dominant integral weights for $\op{SL}(n)$ and $\ell>0 $. Then
\begin{enumerate}
\item For any positive integer $N$, $H^0(\Par_{n,\mn,S},\ml^N)\neq 0$ if and only if
$H^0(\Par_{n,\mn,S},\ml)\neq 0$. Therefore $H^0(\Par_{n,\mn,S},\ml)\neq 0$ if and only if
there is a semi-stable point in $\Par_{n,\mn,S}$ for the polarization given by $\ml$.
\item For any positive integer $N$, $h^0(\Par_{n,\mn,S},\ml^N)=1$ if and only if
$h^0(\Par_{n,\mn,S},\ml)=1$.
\end{enumerate}
\end{proposition}
\begin{proof}
The shift operations reduce (1) to the case $\mn=\mathcal{O}$, this case of (1) follows from
quantum saturation \cite[Theorem 1.8]{BHorn}.

For (2), we can reduce to the case $\mn=\mathcal{O}$, which follows from  \cite[Remark 8.5]{BKq}. Proposition \ref{fulC} below also contains a proof of the quantum Fulton conjecture. See \cite{CS} for a recent discussion of the quantum saturation and Fulton conjectures.
\end{proof}

\begin{proposition}
\source{rigid-local-systems-multiplicative-eigenvalue-problem}
\notready\label{kappa}
\source{rigid-local-systems-multiplicative-eigenvalue-problem}
Suppose $\mathcal{B}(\vec{\lambda},\ell)$ is a line bundle on  $\Par_{n,\mn,S}$ of grade zero where $\lambda^i, i\in[s]$ are dominant integral weights for $\op{SL}(n)$ and $\ell> 0$. Let $D=-\deg(\mn)$. Let $p=p_i\in S$. Define a point $\vec{a}= \Delta_n^s$ as follows: $a^j=\kappa(\frac{\vec{\lambda_j}}{\ell})$  if $j\neq i$ and $a^i=\zeta^D_n\cdot \kappa(\frac{\vec{\lambda_i}}{\ell})$ (see Definition  \ref{symmetry}). Then,
\begin{enumerate}
\item  $\mathcal{B}(\vec{\lambda},\ell)$ is effective if and only $\vec{a}\in P_n(s)$, i.e., there exist matrices $A_i\in \op{SU}(n)$ conjugate to $a^i$, $i=1,\dots,s$ with product $A_1 A_2\cdots A_s=I_n$.
\item  $h^0(\Par_{n,\mn,S},\ml)=1$ if and only if there is a unique $A_i\in \op{SU}(n)$ conjugate to $a^i$, $i=1,\dots,s$  with product $A_1A_2\cdots A_s=I_n$. Here uniqueness means that if $(A'_1,\dots,A'_s)$ is another such collection, there exists a $C\in \op{SU}(n)$ with $A_i'=CA_iC^{-1}$ for $i=1,\dots,s$.
\end{enumerate}
\end{proposition}
\begin{proof}
Suppose $-D=\deg(\mn)\leq 0$. Then $\operatorname{ISh}_p^D: \Par_{n,\mn(D),S}\to \Par_{n,\mn(D),S}$. Let  $\mathcal{B}(\vec{\nu},\ell)\in \Pic^+(\Par_{n,\mn(D),S})$ be the pullback of $\mathcal{B}(\vec{\lambda},\ell)$.  Since $\mn(D)$ is trivial and of grade zero, (1) follows from Proposition \ref{blacktea}, quantum saturation \cite{BHorn}(see Proposition \ref{qTheorems}(1)), and Remark \ref{able}. The case $D<0$ is handled similarly using  $\operatorname{Sh}_p$ instead of  $\operatorname{ISh}_p$.

For (2), we similarly reduce to $D=0$. In this case if $h^0(\Par_{n,\mn,S},\ml)=1$ then by quantum Fulton (Theorem \ref{qTheorems}(2)), $h^0(\Par_{n,\mn,S},\ml^N)=1$ for all $N$ and hence the Mehta-Seshadri moduli space is a point. This shows the uniqueness of the $A_i$. To go the other way, note that  $h^0(\Par_{n,\mn,S},\ml)$ is the space of global sections of  a line bundle over a point and is hence one dimensional.
\end{proof}

The $\kappa$ operation introduces some unnecessary denominators. In fact working in the unitary group $U(n)$, there is an equivalent formulation where the local monodromies are $\ell$th roots of unity. Here $\ell$ is the level of the corresponding line bundle:

\begin{corollary}
\source{rigid-local-systems-multiplicative-eigenvalue-problem}
\notready\label{sansD}
\source{rigid-local-systems-multiplicative-eigenvalue-problem}
Suppose  $\ml=\mathcal{B}(\vec{\lambda},\ell)$ is a line bundle of grade zero on  $\Par_{n,\mn,S}$ where $\lambda^i, i\in[s]$ are dominant integral weights for $\op{SL}(n)$ and $\ell>0$. Write  $\sum_{i=1}^s|\lambda_i|+ \ell \deg\mn=nu$, $u\in \Bbb{Z}$, Then, $\mathcal{B}(\vec{\lambda},\ell)$
is effective on   $\Par_{n,\mn,S}$  if and only if there exist matrices $A_1,\dots,A_s,C\in \op{U}(n)$ with $A_1\cdots A_s=C$, where
\begin{enumerate}
\item $A_i$ is conjugate to the diagonal matrix with entries $\exp(2\pi\sqrt{-1}\lambda^i_j/\ell)$ for $j=1,\dots,n$.
\item $C$ is the central matrix $c I_n$ with $c=\exp(2\pi\sqrt{-1}u/\ell)$.
\end{enumerate}
Moreover,  $h^0(\Par_{n,\mn,S},\ml)=1$ if and only if the corresponding representation variety of unitary local systems  on $\Bbb{P}^1-(S\cup\{p_{s+1}\})$, $p_{s+1}\not\in S$ (with local monodromies $A_1,\dots,A_s, C^{-1}$) is a point (this implies that there is a unique (possibly block reducible) solution to the matrix equation above in $U(n)$ up to conjugation).
\end{corollary}
\begin{proof}
By Proposition \ref{kappa}, $\mathcal{B}(\vec{\lambda},\ell)$
is effective on   $\Par_{n,\mathcal{O},S}$  if and only if there exist matrices $B_1,\dots,B_s\in \op{SU}(n)$ with $B_1\cdots B_s=\zeta_n^{\deg \mn}I_n$, where $A_i$ is conjugate to the diagonal matrix with entries $\exp(2\pi\sqrt{-1}\bigl(\lambda^i_j/{\ell}-\sum |\lambda^i|/(n\ell)\bigr))$ but
$\sum |\lambda^i|/(n\ell)=u/\ell - \deg\mn/n$. This proves the desired statement.

\end{proof}


\subsection{F-vertices and F-line bundles}
\begin{lemma}
\source{rigid-local-systems-multiplicative-eigenvalue-problem}
\notready\label{elgar}
\source{rigid-local-systems-multiplicative-eigenvalue-problem}
Suppose $\ml=\mathcal{O}(E)$  is an F-line bundle on $\Par_{n,\mathcal{O},S}$. Then
\begin{enumerate}
\item $\ml$ gives an extremal ray $\Bbb{Q}_{\geq 0}\mathcal{O}(E)$   of $\Pic^+_{\Bbb{Q}}(\Par_{n,\mathcal{O},S})$, which  lies on some regular face of  $\Pic^+_{\Bbb{Q}}(\Par_{n,\mathcal{O},S})$.
\item $\ml$ gives a Hilbert basis element of the additive semigroup $\Pic^+(\Par_{n,\mathcal{O},S})$.
\end{enumerate}
 \end{lemma}
\begin{proof}
The extremal ray property in (1)  follows from  the method of  \cite[Lemma 2.1]{BHermit}: Suppose
$\ml^m$ is a tensor of two effective line bundles for some positive integer $m$. So we may write
$\ml^m=\mathcal{O}(D_1+D_2)$ where $D_1$ and $D_2$ are effective divisors on the smooth stack $\Par_{n,\mathcal{O},S}$. Since $\ml^m=\mathcal{O}(mE)$, and $\ml^m$ has only one global section up to scalars, we should have $mE=D_1+D_2$ as effective divisors. Since $E$ is irreducible, we see that $D_1$ and $D_2$ are multiples of $E$, and hence the extremal ray property follows. By Proposition \ref{b9}(2), and Lemma \ref{pinky}, $\ml$ lies on some regular face of
$\Pic^+_{\Bbb{Q}}(\Par_{n,\mathcal{O},S})$. Since $\ml$ is an F-line bundle, it is indivisible as an effective line bundle, and so (2) follows from (1).

\end{proof}

\begin{lemma}
\source{rigid-local-systems-multiplicative-eigenvalue-problem}
\notready\label{corresp}
\source{rigid-local-systems-multiplicative-eigenvalue-problem}
There is a bijective correspondence between F-line bundles on $\Par_{n,\mathcal{O},S}$, and F-vertices in $P_n(s)$. The correspondence takes $\mathcal{B}(\vec{\lambda},\ell)$ to $\vec{a}=(\kappa(\lambda^1/\ell),\dots,\kappa(\lambda^s/\ell))$. For the reverse correspondence a vertex $\vec{a}=(a^1,\dots,a^s)$ corresponds to $\mathcal{B}(\vec{\lambda},\ell)$ so that
\begin{enumerate}
\item  $(\kappa(\lambda^1/\ell),\dots,\kappa(\lambda^s/\ell))=\vec{a}$.
\item $\mathcal{B}(\vec{\lambda},\ell)$ is of grade zero (i.e., $\sum |\lambda^i|$ divisible by $n$).
\item The $\ell>0$ above is the smallest possible.
\end{enumerate}
(The condition on $\mathcal{B}(\vec{\lambda},\ell)$ is that it is indivisible as a grade zero line bundle.)
\end{lemma}
\begin{proof}
Let $\ml=(\vec{\lambda},\ell)$ and $\vec{a}$ satisfy the relation $\vec{a}=(\kappa(\lambda^1/\ell),\dots,\kappa(\lambda^s/\ell))$. Then the representation variety of unitary representations with local monodromies given by $\vec{a}$ is identified with the Proj of the graded ring $\oplus_{N=0}^{\infty} H^0(\Par_{n,\mathcal{O},S}, \mathcal{B}(\vec{\lambda},\ell)^N)$. If $\mathcal{B}(\vec{\lambda},\ell)$
is an F-line bundle, then this Proj is a point.  By Lemma \ref{elgar} (1), $\ml$ gives an extremal ray  of $\Pic^+_{\Bbb{Q}}(\Par_{n,\mathcal{O},S})$. This shows the forward direction of the correspondence (also use Proposition \ref{b9}).

For the reverse direction, note that the $\mathcal{B}(\vec{\lambda},\ell)$ described in the statement is effective by quantum saturation (Proposition \ref{qTheorems}). Since (by construction) it is not a multiple of an effective line bundle, the reverse direction follows from Proposition \ref{b9}.
\end{proof}
\begin{remark}
\source{rigid-local-systems-multiplicative-eigenvalue-problem}
\notready\label{inpractice}
\source{rigid-local-systems-multiplicative-eigenvalue-problem}
Note that ``in practice'', to find an F-line bundle corresponding to a vertex $\vec{a}$ we write $a^j=b^j/\tilde{\ell}$ for a common denominator $\tilde{\ell}$, so that $b^j\in \Bbb{Z}^n$ for all $j\in[s]$. Now let $\nu^j=b^j-b^j_n(1,\dots,1)$.
The line bundle $\mathcal{B}(\vec{\nu},\tilde{\ell})$ has grade zero and $(\kappa(\nu^1/\tilde{\ell}),\dots,\kappa(\nu^s/\tilde{\ell}))=\vec{a}$. But it may not be indivisible by this property. We then divide $(\vec{\nu},\tilde{\ell})$ by a positive integer $k$ so that it is indivisible of grade zero. The desired pair $(\lambda,\ell)$ is then $\frac{1}{k}(\vec{\nu},\tilde{\ell})$.
\end{remark}

\begin{lemma}
\source{rigid-local-systems-multiplicative-eigenvalue-problem}
\notready\label{divisorE}
\source{rigid-local-systems-multiplicative-eigenvalue-problem}
Let $\ml=\mathcal{B}(\vec{\lambda},\ell)$ be an F-line bundle on $\Par_{n,\mathcal{O},S}$. Write it as $\mathcal{O}(E)$ for a reduced irreducible divisor  $E\subseteq \Par_{n,\mathcal{O},S}$.
Then $E$ is equal to an effective cycle $C(d,r,\mathcal{O},n,\vec{J})$ (Definition  \ref{cyclist})  of codimension one i.e., $dn  +r(n-r)-\sum_{i=1}^s |\sigma_{J^i}|=1$, with the further property that
$\sum_{i=1}^s \sum_{k\in J^i} \lambda^i_k> d\ell +\sum_{i=1}^s \frac{r|\lambda^i|}{n}$.
\end{lemma}
\begin{proof}
Any (generic) point $(\mw,\me,\gamma)$ of $E$ is unstable for $\ml$. Here we have used the quantum Fulton conjecture (Proposition \ref{qTheorems} (2)) that $h^0(\Par_{n,{\mathcal{O}},S},\ml^N)=1$ for all $N>1$, and hence all sections of $H^0(\Par_{n,{\mathcal{O}},S},\ml^N)$ vanish on $E$.

Let the Harder-Narasimhan maximal contradictor of semistability at a (generic) point $(\mw,\me,\gamma)$ of $E$ be a subbundle $\mv$ of the universal bundle $\mw$ on  $\Par_{n,{\mathcal{O}},S}$ of rank $r$, and suppose it meets the flags of the universal bundle in points corresponding to the Schubert cells parametrized by subsets $J^i\subseteq[n]$, $|J^i|=r$. That is, assume  $\mv_p\in \Omega^0_{J^i}(\me_p)\subseteq \Gr(r,\mw_p)$ for all $p=p_i\in S$.  Let $d=-\deg(\mv)$.

%Note $r>1$ because if $r=1$ then $\ell=1$ by Remark \ref{ajabaja}.
It is now easy to see that $E=C(d,r,\mathcal{O},n,\vec{J})$, and the  codimension of this class is one, i.e., $dn+r(n-r)-\sum_{i=1}^s |\sigma_{J^i}|=1$. This is because the stack $\Omega(d,r,\mathcal{O},n,\vec{J})$ is irreducible and its image under $\pi:\Omega(d,r,\mathcal{O},n,\vec{J})\to
\Par_{n,{\mathcal{O}},S}$ is destabilizing for $\ml$, and hence $\pi$ maps $\Omega(d,r,\mathcal{O},n,\vec{J})$ into $E$. The uniqueness of Harder-Narasimhan maximal contradictors of semistability then shows that $\pi:\Omega(d,r,\mathcal{O},n,\vec{J})\to E$ is generically one-to-one. Therefore $E$ is the cycle theoretic pushforward of $\Omega(d,r,\mathcal{O},n,\vec{J})$. The last inequality is the violated semistability inequality.
\end{proof}
The following gives a more practical approach towards listing all F-line bundles (hence F-vertices of $P_n(s)$ by Lemma \ref{corresp}) on $\Par_{n,\mathcal{O},S}$ in that we do not have to solve systems of linear inequalities.
\begin{proposition}
\source{rigid-local-systems-multiplicative-eigenvalue-problem}
\notready\label{practical}
\source{rigid-local-systems-multiplicative-eigenvalue-problem}
Consider a line bundle  $\ml=\mathcal{B}(\vec{\lambda},\ell)$ which is indivisible as a grade zero line bundle on $\Par_{n,\mathcal{O},S}$. Then $\ml$  is an F-line bundle  iff there is a cycle $E=C(d,r,\mathcal{O},n,\vec{J})$ (Definition  \ref{cyclist})  of codimension one i.e., $dn  +r(n-r)-\sum_{i=1}^s |\sigma_{J^i}|=1$ so that the following conditions hold:
\begin{enumerate}
\item[(i)] $\ml=\mathcal{O}(E)$. Note that we have formulas for $\mathcal{O}(E)$ from Proposition  \ref{enumerative}.
\item[(ii)] $\sum_{i=1}^s \sum_{k\in J^i} \lambda^i_k> d\ell +\sum_{i=1}^s \frac{r|\lambda^i|}{n}.$
\end{enumerate}
Therefore one can list all F-line bundles on $\Par_{n,\mathcal{O},S}$ by listing all cycles $E=C(d,r,\mathcal{O},n,\vec{J})$, computing $\mathcal{O}(E)=\mathcal{B}(\vec{\lambda},\ell)$ from Proposition  \ref{enumerative}, checking if $\mathcal{O}(E)$ satisfies the strict inequality in (ii), and is indivisible as a grade zero line bundle.
\end{proposition}
\begin{proof}
If $\ml=\mathcal{B}(\vec{\lambda},\ell)$ is an F-line bundle then the cycle $C(d,r,\mathcal{O},n,\vec{J})$ exists by the proof of Proposition  \ref{divisorE}. For the reverse implication note that $\mathcal{O}(E)$ is non-semistable on the support of the cycle $C(d,r,\mathcal{O},n,\vec{J})$ (the support is irreducible because it is the image of the irreducible $\Omega(d,r,\mathcal{O},n,\vec{J})$) by (ii), hence any global section of $\ml$ vanishes on this support. The desired assertion follows since  $E$ is not a multiple of an effective divisor by the indivisibility property.
\end{proof}
